% LoRA Architecture Diagram Optimized
\documentclass[tikz,border=15pt]{standalone}
\usepackage{amsmath,amssymb}

% --- UNIVERSAL PREAMBLE BLOCK ---
\usepackage{fontspec}
\usepackage[english, bidi=basic, provide=*]{babel}
\babelprovide[import, onchar=ids fonts]{english}
% 使用通用的 Noto Sans 字体以确保跨平台完美编译
\babelfont{rm}{Noto Sans}
% --------------------------------

% 引入了 decorations.pathreplacing 用于绘制标准的大括号
\usetikzlibrary{shapes,arrows.meta,positioning,calc,backgrounds,decorations.pathreplacing}

\begin{document}

\begin{tikzpicture}[
    node distance=6mm and 8mm,
    % 基础框样式
    box/.style={
        draw=black, thick, rounded corners=3pt,
        minimum width=5.5cm, minimum height=0.8cm, align=center, fill=white,
        font=\normalsize
    },
    % 冻结权重样式（灰色）
    frozen/.style={
        draw=black, thick, rounded corners=3pt,
        fill=black!10, align=center
    },
    % 可训练权重样式（浅蓝色，虚线）
    trainable/.style={
        draw=black, thick, dashed, rounded corners=3pt,
        fill=blue!5, align=center
    },
    % LoRA 标识样式
    lora_badge/.style={
        draw=black, dashed, thick, rounded corners=3pt,
        fill=blue!5, align=center, font=\small, inner sep=4pt
    },
    % 箭头样式
    arrow/.style={
        ->, >=Stealth, thick, black
    },
    % 加法操作符样式
    operator/.style={
        draw=black, thick, circle, minimum size=7mm, inner sep=0pt, align=center, font=\Large
    }
]

% ==========================================
% LEFT: Transformer Layer with LoRA Adapters
% ==========================================
\begin{scope}[local bounding box=left_panel]
    % 节点定义
    \node[box] (input) {\textbf{Input Embeddings}};
    
    \node[box, below=of input] (ln1) {\textbf{Layer Norm}};
    
    \node[box, below=of ln1, minimum height=1.4cm] (mha) {\textbf{Multi-Head Self-Attention}\\ \small (Q, K, V, O projections)};
    
    \node[lora_badge, right=4mm of mha] (lora1) {\textbf{LoRA}\\ \small $r=16$};
    
    \node[box, below=of mha] (add1) {\textbf{Add \& Residual}};
    
    \node[box, below=of add1] (ln2) {\textbf{Layer Norm}};
    
    \node[box, below=of ln2, minimum height=1.2cm] (ffn) {\textbf{FFN}\\ \small (Gate, Up, Down)};
    
    \node[lora_badge, right=4mm of ffn] (lora2) {\textbf{LoRA}\\ \small $r=16$};
    
    \node[box, below=of ffn] (add2) {\textbf{Add \& Residual}};
    
    \node[below=5mm of add2, font=\small\itshape, text=black!60] (output) {Output to Next Layer};

    % 垂直连线
    \draw[arrow] (input) -- (ln1);
    \draw[arrow] (ln1) -- (mha);
    \draw[arrow] (mha) -- (add1);
    \draw[arrow] (add1) -- (ln2);
    \draw[arrow] (ln2) -- (ffn);
    \draw[arrow] (ffn) -- (add2);
    \draw[arrow] (add2) -- (output);

    % LoRA 虚线连接
    \draw[thick, dashed, gray] (mha.east) -- (lora1.west);
    \draw[thick, dashed, gray] (ffn.east) -- (lora2.west);

    % 残差连接 1 (Bypass LN1 & MHA)
    \coordinate (res1) at ($(input.south)!0.5!(ln1.north)$);
    \draw[arrow] (res1) -- ++(-3.5,0) |- (add1.west);

    % 残差连接 2 (Bypass LN2 & FFN)
    \coordinate (res2) at ($(add1.south)!0.5!(ln2.north)$);
    \draw[arrow] (res2) -- ++(-3.5,0) |- (add2.west);

    % --- X24 LAYERS 专业大括号标注 ---
    % 计算右侧括号的确切横坐标，避免与网络重叠
    \coordinate (brace_x) at ($(lora1.east) + (0.4, 0)$);
    \coordinate (brace_top) at (brace_x |- ln1.north);
    \coordinate (brace_bottom) at (brace_x |- add2.south);
    
    % 移除了 mirror 属性，让大括号尖端正确指向右侧文字
    \draw [decorate, decoration={brace, amplitude=8pt}, thick, black!60] 
          (brace_top) -- (brace_bottom) 
          node [midway, right=12pt, font=\large\itshape, text=black!60] {$\times 24$ Layers};

    % 提取左侧标题的绝对 Y 坐标（用于让右侧标题完美对其齐）
    \coordinate (title_y) at ($(input.north) + (0, 0.8)$);
    \node[font=\bfseries\large] at (input.north |- title_y) {Transformer Layer with LoRA Adapters};
    
\end{scope}

% ==========================================
% RIGHT: LoRA Low-Rank Adaptation
% ==========================================
% 大幅向右移动 scope，彻底杜绝文字交叉重叠现象
\begin{scope}[shift={(13.5, 0)}, local bounding box=right_panel]
    
    % 输入节点 (对齐左侧的 input)
    \node[font=\Large, minimum height=0.8cm] (xnode) at (0, 0 |- input) {$x$};
    
    % 连线分叉点（形成对称树状图）
    \coordinate (split) at ($(xnode.south) + (0, -0.4)$);

    % W0 (左侧对称位置)
    \coordinate (w0_center) at (-2.8, -2.8);
    \node[frozen, minimum width=3.5cm, minimum height=2.2cm] (w0) at (w0_center) {\Large $W_0$ \\ \small Frozen pre-trained \\[-2pt] \small weight matrix \\ \small $\in \mathbb{R}^{d \times k}$};

    % A 矩阵 (右侧对称位置，降维)
    \coordinate (a_center) at (2.8, -2.0);
    \node[trainable, minimum width=3.2cm, minimum height=1cm] (A) at (a_center) {\large $A$ \\ \footnotesize $\in \mathbb{R}^{r \times k}$ (Trainable)};

    % B 矩阵 (右侧下方，升维)
    \coordinate (b_center) at (2.8, -3.6);
    \node[trainable, minimum width=3.2cm, minimum height=1cm] (B) at (b_center) {\large $B$ \\ \footnotesize $\in \mathbb{R}^{d \times r}$ (Trainable)};

    % 加法器 (绝对居中)
    \coordinate (sum_center) at (0, -5.5);
    \node[operator] (sum) at (sum_center) {$+$};

    % 输出节点 (绝对居中)
    \coordinate (hout_center) at (0, -7.2);
    \node[box, minimum width=5cm] (hout) at (hout_center) {\large $h = W_0x + \Delta Wx$ \\[-1pt] \normalsize $= W_0x + BAx$};

    % --- 对称路由连线 ---
    % 从输入垂直向下再左右分叉
    \draw[thick] (xnode.south) -- (split);
    \draw[arrow] (split) -| (w0.north);
    \draw[arrow] (split) -| (A.north);

    % A 到 B
    \draw[arrow] (A.south) -- (B.north) node[midway, right, font=\small\itshape] {$r \ll \min(d,k)$};

    % W0 汇聚到 加法器 (标准路径)
    \draw[arrow] (w0.south) |- (sum.west) node[pos=0.75, above, font=\small] {Standard path};
    
    % B 汇聚到 加法器 (LoRA路径)
    \draw[arrow, dashed] (B.south) |- (sum.east) node[pos=0.75, above, font=\small] {\textbf{LoRA path}};

    % 加法器 到 输出
    \draw[arrow] (sum.south) -- (hout.north);

    % 右侧标题 (使用之前保存的 title_y，确保和左边标题百分百水平对齐)
    \node[font=\bfseries\large] at (0, 0 |- title_y) {LoRA Low-Rank Adaptation};

    % 底部说明
    \node[below=6mm of hout, font=\footnotesize\itshape, align=center] {Only $A$ and $B$ are updated during fine-tuning;\\ $W_0$ remains frozen.};
    
\end{scope}

\end{tikzpicture}
\end{document}